%% kapa: exam=1 round=33 year=2022 period=1 form=A subject=real_estate subject_ko=부동산학원론 seq=26 number=106
%% source: Q-Net 감정평가사 기출문제 (cst00309) artlSeq=5209997 file=raw/1008/5209997/2022년도 제33회 감정평가사 1차시험 1교시(민법,경제학,부동산학).pdf
%% conversion: kordoc 4.24.51 -> tools/gen1.py (auto); review=proofread-sonnet; flags=none
\begin{question}{106}
부동산마케팅에 관한 설명으로 옳지 않은 것은?
\begin{choices}
  \item STP란 시장세분화(Segmentation), 표적시장(Target market), 포지셔닝(Positioning)을 말한다.
  \item 마케팅믹스 전략에서의 4P는 유통경로(Place), 제품(Product), 가격(Price), 판매촉진(Promotion)을 말한다.
  \item 노벨티(novelty) 광고는 개인 또는 가정에서 이용되는 실용적이며 장식적인 물건에 상호·전화번호 등을 표시하는 것으로 분양광고에 주로 활용된다.
  \item 관계마케팅 전략은 공급자와 소비자 간의 장기적·지속적인 상호작용을 중요시하는 전략을 말한다.
  \item AIDA 원리에 따르면 소비자의 구매의사결정은 행동(Action), 관심(Interest), 욕망(Desire), 주의(Attention)의 단계를 순차적으로 거친다.
\end{choices}
\end{question}
\answer{5}
